\subsection{对称截断完成化的尺度递归与尾积分证书：\texorpdfstring{$\Xi_{\zeta,\lambda}$}{Xi(zeta,lambda)} 的零点尺度流}\label{subsec:cdim-symmetric-truncation-scale-flow}
Riemann--Siegel--Gabcke 层级以截断阶 \(K\) 在高度变量 \(t\) 上给出可核验的余项界与零点区间递归（\S\ref{subsec:cdim-rs-gabcke-zero-recursion}）。
本小节给出一条互补的“尺度轴”：利用 Jacobi \(\Theta\) 核的模变换把完成化对象的函数方程编译为一族\emph{对称截断} \(\Xi_{\zeta,\lambda}\)；其误差项为显式尾积分，且尺度参数 \(\lambda\) 满足严格的乘法递归与生成方程，从而可将“分辨率”制度化为可计算的重整化流。

\paragraph{Jacobi 核、完成化与对称截断}
令 Jacobi \(\Theta\) 核与完成化 \(\Xi_\zeta\) 分别为
\begin{equation}\label{eq:cdim-jacobi-theta-and-xi}
\boxed{\;
\Theta(x):=\sum_{n\in\ZZ}e^{-\pi n^2x}\quad(x>0),\qquad
\Lambda_\zeta(s):=\pi^{-s/2}\Gamma\!\Bigl(\frac{s}{2}\Bigr)\zeta(s),\qquad
\Xi_\zeta(s):=\frac12 s(s-1)\Lambda_\zeta(s).\;}
\end{equation}
对任意 \(\lambda\ge 1\) 与 \(s\in\CC\setminus\{0,1\}\)，定义对称截断完成化
\begin{equation}\label{eq:cdim-symmetric-truncation-Lambda}
\boxed{\;
\Lambda_{\zeta,\lambda}(s)
:=\frac14\int_{\lambda^{-1}}^{\lambda}\bigl(\Theta(x)-1\bigr)
\Bigl(x^{\frac{s}{2}-1}+x^{\frac{1-s}{2}-1}\Bigr)\,\dd x
\;+\;\frac12\left(
\frac{\lambda^{\frac{s-1}{2}}+\lambda^{-\frac{s-1}{2}}}{s-1}
-\frac{\lambda^{\frac{s}{2}}+\lambda^{-\frac{s}{2}}}{s}
\right),\;}
\end{equation}
并令
\begin{equation}\label{eq:cdim-symmetric-truncation-Xi}
\boxed{\;
\Xi_{\zeta,\lambda}(s):=\frac12 s(s-1)\Lambda_{\zeta,\lambda}(s).\;}
\end{equation}
由于 \eqref{eq:cdim-symmetric-truncation-Lambda} 的极点被 \(s(s-1)\) 抵消，\(\Xi_{\zeta,\lambda}\) 以可去奇点方式延拓为整函数。

\begin{theorem}[对称截断的尺度递归恒等式与尾积分误差项]\label{thm:cdim-symmetric-truncation-tail-identity}
\leanverified{paper\_cdim\_symmetric\_truncation\_tail\_identity}
对任意 \(\lambda\ge 1\) 与 \(s\in\CC\setminus\{0,1\}\)，成立精确分解
\begin{equation}\label{eq:cdim-symmetric-truncation-tail-identity-Lambda}
\boxed{\;
\Lambda_\zeta(s)
=\Lambda_{\zeta,\lambda}(s)
\;+\;\frac12\int_{\lambda}^{\infty}\bigl(\Theta(x)-1\bigr)
\Bigl(x^{\frac{s}{2}-1}+x^{\frac{1-s}{2}-1}\Bigr)\,\dd x.\;}
\end{equation}
从而
\begin{equation}\label{eq:cdim-symmetric-truncation-tail-identity-Xi}
\boxed{\;
\Xi_\zeta(s)
=\Xi_{\zeta,\lambda}(s)
\;+\;\frac14 s(s-1)\int_{\lambda}^{\infty}\bigl(\Theta(x)-1\bigr)
\Bigl(x^{\frac{s}{2}-1}+x^{\frac{1-s}{2}-1}\Bigr)\,\dd x.\;}
\end{equation}
并且 \(\Lambda_{\zeta,\lambda},\Xi_{\zeta,\lambda}\) 满足与原函数一致的自对偶对称性
\begin{equation}\label{eq:cdim-symmetric-truncation-self-duality}
\boxed{\;
\Lambda_{\zeta,\lambda}(s)=\Lambda_{\zeta,\lambda}(1-s),
\qquad
\Xi_{\zeta,\lambda}(s)=\Xi_{\zeta,\lambda}(1-s).\;}
\end{equation}
\end{theorem}
\begin{proof}[证明要点]
当 \(\Re(s)>1\) 时有经典积分表示（参见 \cite{PerezMarco2018NotesRH}）
\[
\Lambda_\zeta(s)=\frac12\int_0^\infty\bigl(\Theta(x)-1\bigr)x^{\frac{s}{2}-1}\,\dd x,
\]
并可将积分区间拆分为 \((0,\lambda^{-1})\cup(\lambda^{-1},\lambda)\cup(\lambda,\infty)\)。
对 \((0,\lambda^{-1})\) 部分作变量代换 \(x=1/u\)，并利用 Jacobi 模变换
\(
\Theta(1/u)=u^{1/2}\Theta(u)
\)
将其改写为 \((\lambda,\infty)\) 上的镜像项与两项初等幂函数校正；再与 \(s\mapsto 1-s\) 的表达式取平均，即得到 \(\Lambda_{\zeta,\lambda}(s)\) 的定义与尾积分项，从而导出 \eqref{eq:cdim-symmetric-truncation-tail-identity-Lambda}。
两侧同为 \(s\) 的亚纯函数且在 \(\Re(s)>1\) 区域一致，解析延拓给出对 \(s\in\CC\setminus\{0,1\}\) 的恒等式。
式 \eqref{eq:cdim-symmetric-truncation-tail-identity-Xi} 由两侧同乘 \(\tfrac12 s(s-1)\) 得出。
自对偶对称性 \eqref{eq:cdim-symmetric-truncation-self-duality} 由 \eqref{eq:cdim-symmetric-truncation-Lambda} 的对称核与校正项在 \(s\mapsto 1-s\) 下不变直接得到。证毕。
\end{proof}

\begin{corollary}[乘法尺度递归律与生成算子]\label{cor:cdim-symmetric-truncation-multiplicative-recursion-generator}
\leanverified{paper\_circle\_dimension\_symmetric\_truncation\_multiplicative\_recursion\_generator}
在定理 \ref{thm:cdim-symmetric-truncation-tail-identity} 的设定下，对任意 \(q>1\)、\(\lambda\ge 1\) 与 \(s\in\CC\setminus\{0,1\}\) 有精确乘法尺度递归
\begin{equation}\label{eq:cdim-symmetric-truncation-multiplicative-recursion}
\boxed{\;
\Lambda_{\zeta,q\lambda}(s)-\Lambda_{\zeta,\lambda}(s)
=\frac12\int_{\lambda}^{q\lambda}\bigl(\Theta(x)-1\bigr)
\Bigl(x^{\frac{s}{2}-1}+x^{\frac{1-s}{2}-1}\Bigr)\,\dd x,\;}
\end{equation}
从而
\begin{equation}\label{eq:cdim-symmetric-truncation-multiplicative-recursion-xi}
\boxed{\;
\Xi_{\zeta,q\lambda}(s)-\Xi_{\zeta,\lambda}(s)
=\frac14 s(s-1)\int_{\lambda}^{q\lambda}\bigl(\Theta(x)-1\bigr)
\Bigl(x^{\frac{s}{2}-1}+x^{\frac{1-s}{2}-1}\Bigr)\,\dd x.\;}
\end{equation}
令 \(r=\log\lambda\)，则 \(\Lambda_{\zeta,e^r}(s)\) 在 \(r\) 上可微并满足生成方程
\begin{equation}\label{eq:cdim-symmetric-truncation-generator}
\boxed{\;
\frac{\partial}{\partial r}\Lambda_{\zeta,e^r}(s)
=\frac12\bigl(\Theta(e^r)-1\bigr)\Bigl(e^{\frac{rs}{2}}+e^{\frac{r(1-s)}{2}}\Bigr),\;}
\end{equation}
以及
\begin{equation}\label{eq:cdim-symmetric-truncation-generator-xi}
\boxed{\;
\frac{\partial}{\partial r}\Xi_{\zeta,e^r}(s)
=\frac14 s(s-1)\bigl(\Theta(e^r)-1\bigr)\Bigl(e^{\frac{rs}{2}}+e^{\frac{r(1-s)}{2}}\Bigr).\;}
\end{equation}
\end{corollary}
\begin{proof}
由 \eqref{eq:cdim-symmetric-truncation-tail-identity-Lambda}，
\(
\Lambda_{\zeta,q\lambda}-\Lambda_{\zeta,\lambda}
=\frac12\int_{\lambda}^{q\lambda}(\Theta(x)-1)\bigl(x^{s/2-1}+x^{(1-s)/2-1}\bigr)\dd x
\)，
即得 \eqref{eq:cdim-symmetric-truncation-multiplicative-recursion}；同乘 \(\tfrac12 s(s-1)\) 得 \eqref{eq:cdim-symmetric-truncation-multiplicative-recursion-xi}。
令 \(q=e^h\) 并令 \(h\to 0\)，由微积分基本定理得到 \eqref{eq:cdim-symmetric-truncation-generator}；
\eqref{eq:cdim-symmetric-truncation-generator-xi} 由定义 \eqref{eq:cdim-symmetric-truncation-Xi} 立即推出。证毕。
\end{proof}

\paragraph{临界线截面与余弦驱动项}
为在纯相位（unitary slice）口径下显式分离尺度参数，定义临界线截面
\[
E_{\zeta,\lambda}(t):=\Xi_{\zeta,\lambda}\!\left(\tfrac12+\mathrm{i}t\right)\in\RR,
\qquad t\in\RR,\ \lambda\ge 1,
\]
并引入正尾核
\[
\psi(x):=\sum_{n\ge 1}e^{-\pi n^2 x}=\frac12\bigl(\Theta(x)-1\bigr)\qquad(x\ge 1).
\]

\begin{proposition}[临界线截断的对数尺度流与余弦驱动项]\label{prop:cdim-symmetric-truncation-critical-line-flow}
\leanverified{paper\_cdim\_symmetric\_truncation\_critical\_line\_flow}
对任意 \(r\ge 0\) 与 \(t\in\RR\) 有
\begin{equation}\label{eq:cdim-symmetric-truncation-critical-line-flow}
\boxed{\;
\frac{\partial}{\partial r}E_{\zeta,e^r}(t)
=-(\tfrac14+t^2)\,\psi(e^r)\,e^{r/4}\cos\!\Bigl(\tfrac{tr}{2}\Bigr)\;}
\end{equation}
并且 \(E_{\zeta,1}(t)\equiv \tfrac12\)，且满足严格积分表示
\begin{equation}\label{eq:cdim-symmetric-truncation-critical-line-integral}
\boxed{\;
E_\zeta(t):=\Xi_\zeta\!\left(\tfrac12+\mathrm{i}t\right)
=\frac12-\int_0^\infty (\tfrac14+t^2)\,\psi(e^r)\,e^{r/4}\cos\!\Bigl(\tfrac{tr}{2}\Bigr)\,\dd r\;}
\end{equation}
\end{proposition}

\begin{proof}
由生成公式 \eqref{eq:cdim-symmetric-truncation-generator-xi}，对 \(s=\tfrac12+\mathrm{i}t\) 有
\[
\frac{\partial}{\partial r}\Xi_{\zeta,e^r}(s)
=\frac14 s(s-1)\bigl(\Theta(e^r)-1\bigr)\Bigl(e^{\frac{rs}{2}}+e^{\frac{r(1-s)}{2}}\Bigr).
\]
注意 \(s(s-1)=-(\tfrac14+t^2)\)，且
\(
e^{\frac{rs}{2}}+e^{\frac{r(1-s)}{2}}
=2e^{r/4}\cos(tr/2)
\)。
代回并用 \(\Theta-1=2\psi\) 即得 \eqref{eq:cdim-symmetric-truncation-critical-line-flow}。
\par
当 \(\lambda=1\) 时，由 \eqref{eq:cdim-symmetric-truncation-Lambda} 可直接化简得
\(\Lambda_{\zeta,1}(s)=\frac{1}{s-1}-\frac{1}{s}\)，从而 \(\Xi_{\zeta,1}(s)\equiv \tfrac12\) 并给出 \(E_{\zeta,1}(t)\equiv\tfrac12\)。
\par
对任意 \(R>0\)，微积分基本定理给出
\(
E_{\zeta,e^R}(t)=\tfrac12+\int_0^R \partial_r E_{\zeta,e^r}(t)\,\dd r
\)。
由定理 \ref{thm:cdim-symmetric-truncation-explicit-error}，当 \(R\to\infty\) 时 \(E_{\zeta,e^R}(t)\to E_\zeta(t)\)，并且积分项绝对可积，故取极限得到 \eqref{eq:cdim-symmetric-truncation-critical-line-integral}。证毕。
\end{proof}

\begin{theorem}[临界带内的显式误差界与尺度深度的指数收敛]\label{thm:cdim-symmetric-truncation-explicit-error}
\leanverified{paper\_cdim\_symmetric\_truncation\_explicit\_error}
对一切 \(\lambda\ge 1\) 与 \(s=\sigma+\mathrm{i}t\)（\(\sigma\in[0,1]\)），成立显式误差界
\begin{equation}\label{eq:cdim-symmetric-truncation-explicit-error-Lambda}
\boxed{\;
\abs{\Lambda_\zeta(s)-\Lambda_{\zeta,\lambda}(s)}
\le
\frac{e^{-\pi\lambda}}{\pi\bigl(1-e^{-3\pi\lambda}\bigr)}
\left(\lambda^{\frac{\sigma}{2}-1}+\lambda^{-\frac{\sigma+1}{2}}\right).\;}
\end{equation}
并且
\begin{equation}\label{eq:cdim-symmetric-truncation-explicit-error-Xi}
\boxed{\;
\abs{\Xi_\zeta(s)-\Xi_{\zeta,\lambda}(s)}
\le
\frac{\abs{s(s-1)}\,e^{-\pi\lambda}}{2\pi\bigl(1-e^{-3\pi\lambda}\bigr)}
\left(\lambda^{\frac{\sigma}{2}-1}+\lambda^{-\frac{\sigma+1}{2}}\right).\;}
\end{equation}
特别地，取 \(\lambda_k=2^k\)（\(k\in\NN\)），则对每个固定 \(s\) 满足 \(\Re(s)\in[0,1]\) 有
\[
\Xi_{\zeta,\lambda_k}(s)\to \Xi_\zeta(s),
\qquad
\abs{\Xi_\zeta(s)-\Xi_{\zeta,\lambda_k}(s)}\ll_s \exp(-\pi 2^k),
\]
从而尺度深度 \(k\) 每推进一层，尾部误差的指数底数即倍增。
\end{theorem}
\begin{proof}
对 \(x\ge \lambda\ge 1\)，由
\(
\Theta(x)-1=2\sum_{n\ge 1}e^{-\pi n^2x}
\)
及不等式 \(n^2\ge 1+3(n-1)\)（\(n\ge 1\)）得
\[
\Theta(x)-1\le 2e^{-\pi x}\sum_{m\ge 0}e^{-3\pi mx}
\le \frac{2e^{-\pi x}}{1-e^{-3\pi\lambda}}.
\]
代入 \eqref{eq:cdim-symmetric-truncation-tail-identity-Lambda} 的尾积分，并注意
\(
\abs{x^{\mathrm{i}t/2}}=1
\)
得到
\[
\abs{\Lambda_\zeta(s)-\Lambda_{\zeta,\lambda}(s)}
\le \frac{1}{1-e^{-3\pi\lambda}}
\int_\lambda^\infty e^{-\pi x}
\left(x^{\frac{\sigma}{2}-1}+x^{-\frac{\sigma+1}{2}}\right)\dd x.
\]
由于 \(x^{\alpha-1}\) 在 \(x\ge 1\) 上单调递减（\(0\le\alpha\le 1\)），有
\(
\int_\lambda^\infty e^{-\pi x}x^{\alpha-1}\dd x
\le \lambda^{\alpha-1}\int_\lambda^\infty e^{-\pi x}\dd x
=e^{-\pi\lambda}\lambda^{\alpha-1}/\pi
\)。
分别取 \(\alpha=\sigma/2\) 与 \(\alpha=(1-\sigma)/2\) 即得 \eqref{eq:cdim-symmetric-truncation-explicit-error-Lambda}；
再同乘 \(\tfrac12\abs{s(s-1)}\) 得 \eqref{eq:cdim-symmetric-truncation-explicit-error-Xi}。证毕。
\end{proof}

\begin{corollary}[有限高度的多项式精度与对数—对数尺度深度]\label{cor:cdim-symmetric-truncation-loglog-depth}
\leanverified{paper\_cdim\_symmetric\_truncation\_loglog\_depth}
给定 \(A>0\) 与 \(T\ge 3\)。若 \(\lambda\ge 1\) 满足
\[
\pi\lambda\ \ge\ (A+3)\log T,
\]
则对所有 \(s=\sigma+\mathrm{i}t\) 且 \(0\le\sigma\le 1,\ |t|\le T\) 有一致估计
\[
\abs{\Xi_\zeta(s)-\Xi_{\zeta,\lambda}(s)}\ \ll_A\ T^{-A}.
\]
因此达到任意多项式阶精度只需 \(\lambda=O_A(\log T)\)，等价地，对数尺度深度 \(r=\log\lambda\) 仅为 \(O_A(\log\log T)\)。
\end{corollary}
\begin{proof}
由定理 \ref{thm:cdim-symmetric-truncation-explicit-error} 的 \eqref{eq:cdim-symmetric-truncation-explicit-error-Xi}，
对 \(0\le\sigma\le 1\) 有
\(
\lambda^{\sigma/2-1}+\lambda^{-(\sigma+1)/2}\le 2\lambda^{-1/2}\le 2
\)；
并且当 \(|t|\le T\) 时，
\(
\abs{s(s-1)}\le (1+t^2)\le 1+T^2
\)。
于是对 \(T\ge 3\) 得到
\[
\sup_{\substack{|t|\le T\\ 0\le\sigma\le 1}}
\abs{\Xi_\zeta(s)-\Xi_{\zeta,\lambda}(s)}
\ \le\ C\,(1+T^2)\,e^{-\pi\lambda}
\ \le\ 2C\,T^2\,e^{-\pi\lambda},
\]
其中 \(C\) 只依赖于 \(1/(1-e^{-3\pi\lambda})\le 1/(1-e^{-3\pi})\) 的上界常数。
若 \(\pi\lambda\ge (A+3)\log T\)，则 \(e^{-\pi\lambda}\le T^{-(A+3)}\)，从而右端 \(\ll_A T^{-A}\)。证毕。
\end{proof}

\begin{theorem}[基于 \texorpdfstring{$\Xi_{\zeta,\lambda}$}{Xi(zeta,lambda)} 的零点计数保持与唯一性认证]\label{thm:cdim-symmetric-truncation-rouche-zero-count}
\leanverified{paper\_cdim\_symmetric\_truncation\_rouche\_zero\_count}
取有界域 \(\Omega\subset\CC\)（分段光滑边界）且 \(\{0,1\}\cap\overline{\Omega}=\varnothing\)。
若存在 \(\lambda\ge 1\) 使
\begin{equation}\label{eq:cdim-symmetric-truncation-rouche-condition}
\sup_{s\in\partial\Omega}\abs{\Xi_\zeta(s)-\Xi_{\zeta,\lambda}(s)}
<
\inf_{s\in\partial\Omega}\abs{\Xi_{\zeta,\lambda}(s)},
\end{equation}
则 \(\Xi_\zeta\) 与 \(\Xi_{\zeta,\lambda}\) 在 \(\Omega\) 内零点个数（按重数计）完全一致。
进一步，若 \(\Xi_{\zeta,\lambda}\) 在 \(\Omega\) 内恰有一个单零点，则 \(\Xi_\zeta\) 在 \(\Omega\) 内亦恰有一个单零点。
\par
并且 \eqref{eq:cdim-symmetric-truncation-rouche-condition} 的左侧可由定理 \ref{thm:cdim-symmetric-truncation-explicit-error} 的显式界替换为可计算的充分条件：
若 \(\partial\Omega\) 落在 \(\Re(s)\in[0,1]\) 内，则
\begin{equation}\label{eq:cdim-symmetric-truncation-rouche-sufficient}
\frac{e^{-\pi\lambda}}{2\pi\bigl(1-e^{-3\pi\lambda}\bigr)}
\sup_{s\in\partial\Omega}\abs{s(s-1)}
\cdot
\sup_{s\in\partial\Omega}
\left(\lambda^{\frac{\Re(s)}{2}-1}+\lambda^{-\frac{\Re(s)+1}{2}}\right)
<
\inf_{s\in\partial\Omega}\abs{\Xi_{\zeta,\lambda}(s)}
\end{equation}
则零点计数与唯一性认证成立。
\end{theorem}
\begin{proof}
条件 \eqref{eq:cdim-symmetric-truncation-rouche-condition} 等价于 \(\abs{\Xi_\zeta-\Xi_{\zeta,\lambda}}<\abs{\Xi_{\zeta,\lambda}}\) 在 \(\partial\Omega\) 上处处成立，
由 Rouch\'e 定理知 \(\Xi_\zeta\) 与 \(\Xi_{\zeta,\lambda}\) 在 \(\Omega\) 内零点数相同。
式 \eqref{eq:cdim-symmetric-truncation-rouche-sufficient} 由定理 \ref{thm:cdim-symmetric-truncation-explicit-error} 的上界对 \(\partial\Omega\) 取上确界直接得到。证毕。
\end{proof}

\begin{theorem}[单零点的尺度流方程：由生成算子诱导的确定性零点递归]\label{thm:cdim-symmetric-truncation-zero-flow}
\leanverified{paper\_cdim\_symmetric\_truncation\_zero\_flow}
取 \(\lambda_0>1\)，并设 \(\rho_0\) 为 \(\Xi_{\zeta,\lambda_0}\) 的单零点，即
\[
\Xi_{\zeta,\lambda_0}(\rho_0)=0,
\qquad
\Xi_{\zeta,\lambda_0}'(\rho_0)\neq 0.
\]
则存在 \(\lambda_0\) 的邻域 \(I\subset(0,\infty)\) 与唯一解析函数 \(\rho:I\to\CC\) 使
\[
\rho(\lambda_0)=\rho_0,
\qquad
\Xi_{\zeta,\lambda}(\rho(\lambda))\equiv 0\quad(\lambda\in I).
\]
令 \(r=\log\lambda\) 并写 \(\rho(r):=\rho(e^r)\)，则 \(\rho\) 满足精确零点流方程
\begin{equation}\label{eq:cdim-symmetric-truncation-zero-flow}
\boxed{\;
\frac{\dd\rho}{\dd r}
=-\frac{\partial_r\Xi_{\zeta,e^r}(\rho)}{\Xi_{\zeta,e^r}'(\rho)}
=-\frac{\frac14\,\rho(\rho-1)\bigl(\Theta(e^r)-1\bigr)\Bigl(e^{\frac{r\rho}{2}}+e^{\frac{r(1-\rho)}{2}}\Bigr)}
{\Xi_{\zeta,e^r}'(\rho)}.\;}
\end{equation}
因此，对任意小步长 \(h\) 有一阶递归更新
\begin{equation}\label{eq:cdim-symmetric-truncation-zero-flow-euler}
\boxed{\;
\rho(r+h)
=\rho(r)-h\,\frac{\partial_r\Xi_{\zeta,e^r}(\rho(r))}{\Xi_{\zeta,e^r}'(\rho(r))}
+O(h^2)\qquad(h\to 0).\;}
\end{equation}
\end{theorem}
\begin{proof}
由 \((\lambda,s)\mapsto \Xi_{\zeta,\lambda}(s)\) 的解析性与 \(\Xi_{\zeta,\lambda_0}'(\rho_0)\neq 0\)，隐函数定理给出局部唯一解析零点分支 \(\rho(\lambda)\)。
对恒等式 \(\Xi_{\zeta,e^r}(\rho(r))\equiv 0\) 对 \(r\) 求导得到
\[
0=\partial_r\Xi_{\zeta,e^r}(\rho)+\Xi_{\zeta,e^r}'(\rho)\,\frac{\dd\rho}{\dd r},
\]
从而得到第一等式；第二等式由生成公式 \eqref{eq:cdim-symmetric-truncation-generator-xi} 直接代入得到。
式 \eqref{eq:cdim-symmetric-truncation-zero-flow-euler} 为对 \(\rho(r+h)\) 的 Taylor 展开。证毕。
\end{proof}

\begin{corollary}[尾积分误差诱导的残差证书：尺度推进下的零点极限]\label{cor:cdim-symmetric-truncation-zero-branch-residual-certificate}
\leanverified{paper\_cdim\_symmetric\_truncation\_zero\_branch\_residual\_certificate}
在定理 \ref{thm:cdim-symmetric-truncation-zero-flow} 的设定下，
设零点分支可延拓至 \([\lambda_0,\infty)\)，并且沿分支存在常数 \(m>0\) 使
\(
\abs{\Xi_{\zeta,\lambda}'(\rho(\lambda))}\ge m
\)。
若进一步满足 \(\Re(\rho(\lambda))\in[0,1]\)（\(\lambda\ge \lambda_0\)），则对每个 \(\lambda\ge \lambda_0\) 有残差恒等式
\[
\abs{\Xi_\zeta(\rho(\lambda))}
=\abs{\Xi_\zeta(\rho(\lambda))-\Xi_{\zeta,\lambda}(\rho(\lambda))}
\]
并可由 \eqref{eq:cdim-symmetric-truncation-explicit-error-Xi} 给予显式上界。
特别地，若 \(\rho(\lambda)\to\rho_\infty\) 且 \(\Xi_\zeta'(\rho_\infty)\neq 0\)，则 \(\rho_\infty\) 为 \(\Xi_\zeta\) 的单零点，且
\(
\abs{\rho(\lambda)-\rho_\infty}
\)
以与 \(\abs{\Xi_\zeta(\rho(\lambda))}\) 同阶的速度衰减。
\end{corollary}
\begin{proof}[证明要点]
由 \(\Xi_{\zeta,\lambda}(\rho(\lambda))=0\) 得残差恒等式成立。
误差上界由定理 \ref{thm:cdim-symmetric-truncation-explicit-error} 代入 \(s=\rho(\lambda)\) 得出。
若 \(\rho(\lambda)\to\rho_\infty\)，则由 \(\Xi_\zeta(\rho(\lambda))\to 0\) 得 \(\Xi_\zeta(\rho_\infty)=0\)。
当 \(\Xi_\zeta'(\rho_\infty)\neq 0\) 时，\(\Xi_\zeta\) 在 \(\rho_\infty\) 的邻域为双射到其像，故 \(\abs{\rho(\lambda)-\rho_\infty}\asymp \abs{\Xi_\zeta(\rho(\lambda))}\)。证毕。
\end{proof}

\begin{theorem}[临界线零点分支的超指数漂移界]\label{thm:cdim-symmetric-truncation-critical-line-zero-drift}
\leanverified{paper\_cdim\_symmetric\_truncation\_critical\_line\_zero\_drift}
取 \(T>0\) 与 \(r_0\ge 0\)。
假设存在连续函数 \(\gamma:[r_0,\infty)\to[-T,T]\) 使得对所有 \(r\ge r_0\),
\[
E_{\zeta,e^r}(\gamma(r))=0,
\qquad
\partial_t E_{\zeta,e^r}(\gamma(r))\neq 0,
\qquad
\inf_{r\ge r_0}\bigl|\partial_t E_{\zeta,e^r}(\gamma(r))\bigr|\ge m>0.
\]
则极限 \(\gamma(\infty):=\lim_{r\to\infty}\gamma(r)\) 存在，且满足显式位移界
\begin{equation}\label{eq:cdim-symmetric-truncation-critical-line-zero-drift}
\boxed{\;
|\gamma(\infty)-\gamma(r_0)|
\le
\frac{\tfrac14+T^2}{\pi\,m\,(1-e^{-\pi})}\;
e^{-3r_0/4}\;
e^{-\pi e^{r_0}}\;}
\end{equation}
\end{theorem}
\begin{proof}
由隐函数定理，\(\gamma\) 在每个满足 \(\partial_t E_{\zeta,e^r}(\gamma(r))\neq 0\) 的点可微，且
\[
\gamma'(r)=-\frac{\partial_r E_{\zeta,e^r}(\gamma(r))}{\partial_t E_{\zeta,e^r}(\gamma(r))}.
\]
由 \eqref{eq:cdim-symmetric-truncation-critical-line-flow} 与 \(|\cos|\le 1\) 得
\[
\bigl|\partial_r E_{\zeta,e^r}(t)\bigr|
\le (\tfrac14+t^2)\,\psi(e^r)\,e^{r/4}
\le (\tfrac14+T^2)\,\psi(e^r)\,e^{r/4}
\qquad(|t|\le T).
\]
对 \(x\ge 1\)，由 \(\psi(x)=\sum_{n\ge 1}e^{-\pi n^2 x}\le \sum_{n\ge 1}e^{-\pi n x}=\frac{e^{-\pi x}}{1-e^{-\pi x}}\le \frac{e^{-\pi x}}{1-e^{-\pi}}\) 得
\(
\psi(e^r)\le \frac{e^{-\pi e^r}}{1-e^{-\pi}}
\)。
结合 \(|\partial_t E_{\zeta,e^r}(\gamma(r))|\ge m\) 得
\[
|\gamma'(r)|
\le
\frac{\tfrac14+T^2}{m(1-e^{-\pi})}\,e^{r/4}e^{-\pi e^r}.
\]
右端在 \([r_0,\infty)\) 可积，故 \(\gamma\) 为 Cauchy，从而极限 \(\gamma(\infty)\) 存在且
\[
|\gamma(\infty)-\gamma(r_0)|
\le \int_{r_0}^{\infty}|\gamma'(r)|\,\dd r
\le
\frac{\tfrac14+T^2}{m(1-e^{-\pi})}\int_{r_0}^{\infty}e^{r/4}e^{-\pi e^r}\,\dd r.
\]
对尾积分作代换 \(v=\pi e^r\)（于是 \(\dd r=\dd v/v\)）得
\[
\int_{r_0}^{\infty}e^{r/4}e^{-\pi e^r}\,\dd r
=\pi^{-1/4}\int_{\pi e^{r_0}}^{\infty} v^{-3/4}e^{-v}\,\dd v
=\pi^{-1/4}\Gamma\!\left(\tfrac14,\pi e^{r_0}\right).
\]
由 \(0<\tfrac14\le 1\) 时不完全伽马函数的基本界
\(
\Gamma(\alpha,z)\le z^{\alpha-1}e^{-z}
\)
得
\[
\Gamma\!\left(\tfrac14,\pi e^{r_0}\right)\le (\pi e^{r_0})^{-3/4}e^{-\pi e^{r_0}},
\]
代回即得 \eqref{eq:cdim-symmetric-truncation-critical-line-zero-drift}。证毕。
\end{proof}

\begin{proposition}[对称解析形变中的临界线多重零点屏障]\label{prop:cdim-symmetric-truncation-multiple-zero-barrier}
\leanverified{paper\_cdim\_symmetric\_truncation\_multiple\_zero\_barrier}
设 \(\{f_\lambda\}_{\lambda\in I}\) 为参数区间 \(I\subset\RR\) 上的整函数族，且 \((\lambda,s)\mapsto f_\lambda(s)\) 联合解析，并满足对称性
\[
f_\lambda(s)=f_\lambda(1-s)\qquad(\lambda\in I,\ s\in\CC).
\]
若存在连续零点轨道 \(\lambda\mapsto \rho(\lambda)\) 使得对某 \(\lambda_*\in I\) 有 \(\Re\rho(\lambda_*)\neq\tfrac12\)，而对另一个 \(\lambda^*\in I\) 有 \(\Re\rho(\lambda^*)=\tfrac12\)，则存在 \(\lambda^\dagger\in I\) 与 \(t^\dagger\in\RR\) 使
\[
f_{\lambda^\dagger}\!\left(\tfrac12+\mathrm{i}t^\dagger\right)=0,
\qquad
\partial_s f_{\lambda^\dagger}\!\left(\tfrac12+\mathrm{i}t^\dagger\right)=0.
\]
换言之，离轴零点分支若在对称形变中穿越临界线，则必经临界线上的多重零点事件。
\end{proposition}
\begin{proof}
当 \(\Re\rho(\lambda)\neq\tfrac12\) 时，由对称性 \(f_\lambda(s)=f_\lambda(1-s)\) 可知 \(\rho(\lambda)\) 与其镜像 \(1-\rho(\lambda)\) 为不同点且同为零点。
若存在参数使 \(\rho(\lambda)\to \tfrac12+\mathrm{i}t^\dagger\)，则同时 \(1-\rho(\lambda)\to \tfrac12+\mathrm{i}t^\dagger\)，即两条不同零点分支在临界线上汇合。
解析函数零点分支的汇合必对应多重零点（否则由隐函数定理零点分支在该点局部唯一而不能合并），从而得到所述 \(\partial_s f_{\lambda^\dagger}(\tfrac12+\mathrm{i}t^\dagger)=0\)。证毕。
\end{proof}

\begin{conjecture}[尺度流的临界线不变性与吸引性]\label{conj:cdim-symmetric-truncation-critical-line}
设 \(\rho_\lambda\) 为 \(\Xi_{\zeta,\lambda}\) 的非平凡零点。
以下两种陈述任一成立都将推出 RH：
\begin{enumerate}
  \item \textbf{不变性：}对每个 \(\lambda\ge 1\)，\(\Xi_{\zeta,\lambda}\) 的全部非平凡零点满足 \(\Re(\rho_\lambda)=\tfrac12\)。
  \item \textbf{吸引性：}存在 \(\lambda_\star\ge 1\)，使得对所有 \(\lambda\ge \lambda_\star\)，\(\Xi_{\zeta,\lambda}\) 的全部非平凡零点满足 \(\Re(\rho_\lambda)=\tfrac12\)。
\end{enumerate}
\end{conjecture}

\begin{remark}[尺度流猜想 \(\Rightarrow\) RH 的极限传递机制]\label{rem:cdim-symmetric-truncation-critical-line-implies-rh}
由定理 \ref{thm:cdim-symmetric-truncation-explicit-error}，\(\Xi_{\zeta,\lambda}\to \Xi_\zeta\) 在任意紧集上一致收敛。
在该一致收敛与零点不在边界的条件下，Hurwitz 定理将“零点落在临界线”的性质从 \(\Xi_{\zeta,\lambda}\) 传递到极限函数 \(\Xi_\zeta\)。
因此，猜想 \ref{conj:cdim-symmetric-truncation-critical-line} 的任一形式均可视作 RH 的一种动力学表述：临界线成为尺度流下的全局不变集或最终吸引集。
离散 dyadic 壳层版本的交错稳定性猜想（猜想 \ref{conj:cdim-symmetric-truncation-dyadic-lp-interlacing}）提供另一条可核验入口：若足够深度的 dyadic 截断已落入 Laguerre--P\'olya 类并呈现严格交错，则相同的 Hurwitz 机制可把“实零点性/临界线性”传递到极限函数 \(E_\zeta(t)=\Xi_\zeta(\tfrac12+\mathrm{i}t)\)。
\end{remark}

\begin{conjecture}[dyadic 壳层截断的 Laguerre--P\'olya 稳定与零点交错]\label{conj:cdim-symmetric-truncation-dyadic-lp-interlacing}
令 \(\lambda_k=2^k\) 并记
\[
E_k(t):=E_{\zeta,\lambda_k}(t)=\Xi_{\zeta,\lambda_k}\!\left(\tfrac12+\mathrm{i}t\right)
\qquad(k\in\NN,\ t\in\RR).
\]
猜想存在 \(k_0\ge 1\) 使得对所有 \(k\ge k_0\) 成立：
\begin{enumerate}
  \item \(E_k\) 属于 Laguerre--P\'olya 类，因而其全部零点为实且互异；
  \item \(E_{k+1}\) 的零点与 \(E_k\) 的零点严格交错；
  \item 若 \(\gamma_{k,j}\) 为 \(E_k\) 的第 \(j\) 个正零点，则 \(\gamma_{k,j}\to\gamma_j\)（\(k\to\infty\)）且 \(\{\pm\gamma_j\}\) 给出 \(E_\zeta\) 的对应零点；并且存在常数 \(C_j\)（允许多项式增长）使
  \[
  |\gamma_j-\gamma_{k,j}|\ \le\ C_j\,e^{-\pi 2^k}\qquad(k\to\infty).
  \]
\end{enumerate}
\end{conjecture}

\begin{remark}[交错猜想的可核验性]\label{rem:cdim-symmetric-truncation-dyadic-lp-auditability}
定理 \ref{thm:cdim-symmetric-truncation-explicit-error} 已给出 dyadic 深度 \(\lambda_k=2^k\) 下的显式尾项界 \(\abs{E_\zeta(t)-E_k(t)}\ll_{t}e^{-\pi 2^k}\)。
因此，若能在某个 \(k_0\) 之后建立（或数值核验并进一步寻求证明）Laguerre--P\'olya 性与交错性，则由 Hurwitz 稳定性可把实零点性从 \(\{E_k\}\) 传递到 \(E_\zeta\)，从而得到 RH。
\end{remark}

上述对称截断框架把“尺度递归—尾积分误差—零点尺度流”组织为一条可审计链：尺度参数 \(\lambda\) 在乘法递归下累积 \(\Theta\) 核信息，而尾积分在 \(e^{-\pi\lambda}\) 级显式衰减并给出可计算误差证书（定理 \ref{thm:cdim-symmetric-truncation-explicit-error}）。
Rouch\'e 条件提供有限域内零点计数/唯一性认证（定理 \ref{thm:cdim-symmetric-truncation-rouche-zero-count}），生成算子诱导的流方程把零点的尺度演化写成确定性递归（定理 \ref{thm:cdim-symmetric-truncation-zero-flow}），并在临界线分支非退化时给出超指数漂移压制（定理 \ref{thm:cdim-symmetric-truncation-critical-line-zero-drift}）。
对称解析形变下离轴分支若穿越临界线则必经多重零点事件（命题 \ref{prop:cdim-symmetric-truncation-multiple-zero-barrier}），从而把“离轴生成/消失”转写为可定位的退化见证并与审计口径对齐。

\endinput
